\chapter{GCD, the Euclidean Algorithm and Unique Factorization}

\section{Greatest common divisor}

\begin{definition}
Let $a,b \in \Z$, not both zero. The \emph{greatest common divisor} $\gcd(a,b)$ is the largest
integer $d$ with $d \mid a$ and $d \mid b$. We set $\gcd(0,0) = 0$.
\end{definition}

Since $1$ is always a common divisor and (by BBD) every common divisor is at most
$\max(\abs a,\abs b)$, the gcd exists and is positive when $a,b$ are not both $0$.
Simple facts: $\gcd(a,b) = \gcd(b,a) = \gcd(\abs a,\abs b)$, $\gcd(a,0) = \abs a$, and
$\gcd(a,1) = 1$.

\begin{proposition}[GCD With Remainders, GCDWR]
If $a = qb + r$ for integers $a,b,q,r$, then $\gcd(a,b) = \gcd(b,r)$.
\end{proposition}
\begin{proof}
If $d$ divides $a$ and $b$, then by DIC $d \mid a - qb = r$. If $d$ divides $b$ and $r$, then
$d \mid qb + r = a$. So $a,b$ and $b,r$ have exactly the same common divisors, hence the same
greatest one.
\end{proof}

\section{The Euclidean algorithm}

Applying GCDWR repeatedly with the division algorithm shrinks the numbers until a remainder
is $0$; the last non-zero remainder is the gcd.

\begin{example}
Compute $\gcd(1386, 322)$.
\begin{align*}
1386 &= 4\cdot 322 + 98\\
322 &= 3 \cdot 98 + 28\\
98 &= 3\cdot 28 + 14\\
28 &= 2\cdot 14 + 0.
\end{align*}
So $\gcd(1386,322) = \gcd(322,98) = \gcd(98,28) = \gcd(28,14) = \gcd(14,0) = 14$.
\end{example}

The remainders strictly decrease and are non-negative, so the algorithm terminates.

\section{Bezout's lemma and the extended Euclidean algorithm}

\begin{theorem}[Bezout's Lemma / EEA]
For all $a,b \in \Z$ there exist integers $x,y$ with $ax + by = \gcd(a,b)$.
\end{theorem}

The \emph{extended Euclidean algorithm} finds such $x,y$. We maintain rows $(x_i, y_i, r_i)$
with the invariant
\[ a x_i + b y_i = r_i. \]
Start with rows $(1, 0, a)$ and $(0, 1, b)$. Given two consecutive rows, let
$q_{i} = \floor{r_{i-2}/r_{i-1}}$ and set
\[ \text{row}_i = \text{row}_{i-2} - q_i\cdot \text{row}_{i-1}. \]
The invariant is preserved because it is linear. Stop when $r = 0$; the previous row gives the
gcd and the coefficients. (The $r$ column is exactly the Euclidean algorithm.)

\begin{example}[EEA for $a = 1386$, $b = 322$]
\begin{center}
\begin{tabular}{rrrr}
\toprule
$x$ & $y$ & $r$ & $q$ \\
\midrule
$1$ & $0$ & $1386$ & \\
$0$ & $1$ & $322$ & \\
$1$ & $-4$ & $98$ & $4$\\
$-3$ & $13$ & $28$ & $3$\\
$10$ & $-43$ & $14$ & $3$\\
$-23$ & $99$ & $0$ & $2$\\
\bottomrule
\end{tabular}
\end{center}
For instance the fourth row is $(0,1,322) - 3(1,-4,98) = (-3, 13, 28)$. So
$\gcd(1386,322) = 14$ and
\[ 1386(10) + 322(-43) = 13860 - 13846 = 14. \]
(The last row gives $1386(-23) + 322(99) = 0$, which is useful later for generating all
solutions.)
\end{example}

\begin{example}[EEA for $a = 506$, $b = 391$]
\begin{center}
\begin{tabular}{rrrr}
\toprule
$x$ & $y$ & $r$ & $q$ \\
\midrule
$1$ & $0$ & $506$ & \\
$0$ & $1$ & $391$ & \\
$1$ & $-1$ & $115$ & $1$\\
$-3$ & $4$ & $46$ & $3$\\
$7$ & $-9$ & $23$ & $2$\\
$-17$ & $22$ & $0$ & $2$\\
\bottomrule
\end{tabular}
\end{center}
$\gcd(506,391) = 23$ and $506(7) + 391(-9) = 3542 - 3519 = 23$.
\end{example}

\begin{important}[Back-substitution]
Instead of the table, one can run the Euclidean algorithm and substitute backwards:
$14 = 98 - 3\cdot 28 = 98 - 3(322 - 3\cdot 98) = 10\cdot 98 - 3\cdot 322
= 10(1386 - 4\cdot 322) - 3\cdot 322 = 10\cdot 1386 - 43\cdot 322$. Same answer; the table is
less error-prone.
\end{important}

\begin{theorem}[GCD Characterization Theorem, GCDCT]
Let $a,b \in \Z$ and $d \geq 0$. If $d \mid a$, $d \mid b$, and $d = ax + by$ for some integers
$x,y$, then $d = \gcd(a,b)$.
\end{theorem}
\begin{proof}
Let $g = \gcd(a,b)$. By DIC, $g \mid ax + by = d$, so (if $d > 0$) $g \le d$ by BBD. But $d$
is a common divisor, so $d \le g$. Hence $d = g$. (If $d = 0$ then $a = b = 0$.)
\end{proof}

\begin{important}[Bezout does not go backwards]
$ax + by = d$ alone does not make $d$ the gcd: $6\cdot 1 + 4 \cdot 1 = 10$ but
$\gcd(6,4) = 2$. You also need $d$ to be a common divisor. What \emph{is} true is that the
set of integer combinations $\{ax + by\}$ is exactly the set of multiples of $\gcd(a,b)$.
\end{important}

\begin{corollary}[Common Divisor Divides GCD, CDDG]
If $c \mid a$ and $c \mid b$ then $c \mid \gcd(a,b)$.
\end{corollary}
\begin{proof}
Write $\gcd(a,b) = ax + by$ and apply DIC.
\end{proof}

\section{Coprimeness}

\begin{definition}
$a$ and $b$ are \emph{coprime} (relatively prime) if $\gcd(a,b) = 1$.
\end{definition}

\begin{theorem}[Coprimeness Characterization Theorem, CCT]
$\gcd(a,b) = 1$ if and only if there exist integers $x,y$ with $ax + by = 1$.
\end{theorem}
\begin{proof}
($\Rightarrow$) Bezout. ($\Leftarrow$) $1$ divides $a$ and $b$, so by GCDCT $1 = \gcd(a,b)$.
\end{proof}

\begin{proposition}[Division by GCD, DB GCD]
If $d = \gcd(a,b) \ne 0$, then $\gcd(a/d,\, b/d) = 1$.
\end{proposition}
\begin{proof}
From $ax + by = d$, divide by $d$: $(a/d)x + (b/d)y = 1$. Apply CCT.
\end{proof}

\begin{proposition}[Coprimeness and Divisibility, CAD]
If $c \mid ab$ and $\gcd(a,c) = 1$, then $c \mid b$.
\end{proposition}
\begin{proof}
Write $ax + cy = 1$. Multiply by $b$: $abx + cby = b$. Now $c \mid ab$ and $c \mid c$, so by
DIC $c \mid abx + cby = b$.
\end{proof}

\begin{example}
Prove that for any integer $n$, $\gcd(5n+3,\, 3n+2) = 1$: $3(5n+3) - 5(3n+2) = -1$, so
$(5n+3)(-3) + (3n+2)(5) = 1$, and CCT applies.
\end{example}

\section{Primes and Euclid's lemma}

\begin{theorem}[Euclid's Lemma / Primes and Divisibility, PAD]
If $p$ is prime and $p \mid ab$, then $p \mid a$ or $p \mid b$.
\end{theorem}
\begin{proof}
Suppose $p \nmid a$. The only positive divisors of $p$ are $1$ and $p$, and $p \nmid a$, so
$\gcd(p,a) = 1$. By CAD, $p \mid b$.
\end{proof}

By induction, if $p \mid a_1a_2\cdots a_n$ then $p \mid a_i$ for some $i$.

\begin{example}
$\sqrt p$ is irrational for every prime $p$: if $p q^2 = m^2$ with $\gcd(m,q) = 1$, then
$p \mid m^2$ so $p \mid m$ (PAD); write $m = pk$, so $q^2 = pk^2$ and $p \mid q$, contradicting
$\gcd(m,q) = 1$.
\end{example}

\section{Unique factorization}

\begin{theorem}[Unique Factorization Theorem, UFT]
Every integer $n > 1$ can be written as a product of primes, and this factorization is
unique up to the order of the factors.
\end{theorem}
\begin{proof}
Existence was proved by strong induction in Chapter 4. For uniqueness, suppose
$n = p_1p_2\cdots p_k = q_1q_2 \cdots q_\ell$ with all $p_i,q_j$ prime; induct on $k$. Since
$p_1 \mid q_1\cdots q_\ell$, PAD gives $p_1 \mid q_j$ for some $j$, and since $q_j$ is prime,
$p_1 = q_j$. Cancel it from both sides and apply the inductive hypothesis to the shorter
product. (If $k = 1$, $n = p_1$ is prime, which forces $\ell = 1$.)
\end{proof}

We often write the \emph{prime factorization} as $n = p_1^{\alpha_1}p_2^{\alpha_2}\cdots
p_k^{\alpha_k}$ with distinct primes and $\alpha_i \geq 1$; allowing $\alpha_i = 0$ lets us
write two numbers over the same list of primes.

\begin{proposition}[Divisors from Prime Factorization]
If $n = p_1^{\alpha_1}\cdots p_k^{\alpha_k}$, then the positive divisors of $n$ are exactly the
numbers $p_1^{\beta_1}\cdots p_k^{\beta_k}$ with $0 \le \beta_i \le \alpha_i$. So $n$ has
$(\alpha_1 + 1)\cdots(\alpha_k + 1)$ positive divisors.
\end{proposition}

\begin{theorem}[GCD from Prime Factorization, GCDPF]
If $a = \prod p_i^{\alpha_i}$ and $b = \prod p_i^{\beta_i}$ (with $\alpha_i,\beta_i \ge 0$),
then
\[ \gcd(a,b) = \prod p_i^{\min(\alpha_i,\beta_i)}, \qquad
   \lcm(a,b) = \prod p_i^{\max(\alpha_i,\beta_i)}. \]
In particular $\gcd(a,b)\lcm(a,b) = ab$ for positive $a,b$.
\end{theorem}

\begin{example}
$1386 = 2\cdot 3^2 \cdot 7 \cdot 11$ and $322 = 2 \cdot 7 \cdot 23$. Then
$\gcd = 2\cdot 7 = 14$ (agreeing with the Euclidean algorithm) and
$\lcm = 2\cdot 3^2\cdot 7\cdot 11\cdot 23 = 31878 = 1386\cdot 322/14$.
$1386$ has $(1+1)(2+1)(1+1)(1+1) = 24$ positive divisors.
\end{example}

\begin{aside}
Factoring large numbers is believed to be hard, while the Euclidean algorithm is very fast
(the number of steps is at most about five times the number of decimal digits). That contrast
is the basis of RSA (Chapter 10).
\end{aside}

\section{Common pitfalls}
\begin{itemize}
  \item Concluding $d = \gcd(a,b)$ from $ax + by = d$ alone (GCDCT needs $d$ to be a common
        divisor too).
  \item Arithmetic slips in the EEA table; always verify $ax + by = d$ at the end.
  \item Using PAD for a non-prime: $4 \mid 2\cdot 6$ but $4\nmid 2$, $4 \nmid 6$.
\end{itemize}

\section{Summary}
\begin{itemize}
  \item GCDWR: $\gcd(a,b) = \gcd(b,r)$; iterating gives the Euclidean algorithm.
  \item EEA: rows $(x,y,r)$ with $ax + by = r$; produces Bezout coefficients.
  \item GCDCT: a non-negative common divisor that is an integer combination is the gcd.
        CCT: $\gcd(a,b) = 1 \iff ax+by = 1$ is solvable.
  \item CAD: $c\mid ab$, $\gcd(a,c) = 1 \Rightarrow c \mid b$. PAD (Euclid's lemma): $p \mid ab
        \Rightarrow p\mid a$ or $p \mid b$.
  \item UFT: unique prime factorization; gcd/lcm via min/max of exponents.
\end{itemize}
